% Grado 6 -- generado por tools/build_tex.py a partir de raw/grado_06.txt
\section{Grado 6}\label{grado:6}

% PDF p. 45
\dba{1}{Interpreta los números enteros y racionales (en sus representaciones de fracción y de decimal) con sus operaciones, en diferentes contextos, al resolver problemas de variación, repartos, particiones, estimaciones, etc. Reconoce y establece diferentes relaciones (de orden y equivalencia y las utiliza para argumentar procedimientos).}

\evidencias

\begin{itemize}
  \item Resuelve problemas en los que intervienen cantidades positivas y negativas en procesos de comparación, transformación y representación.
  \item Propone y justifica diferentes estrategias para resolver problemas con números enteros, racionales (en sus representaciones de fracción y de decimal) en contextos escolares y extraescolares.
  \item Representa en la recta numérica la posición de un número utilizando diferentes estrategias.
  \item Interpreta y justifica cálculos numéricos al solucionar problemas.
\end{itemize}

\ejemplo

En una competencia de autos a tres vueltas, el tiempo (en minutos) de cada vuelta se registró en la siguiente tabla.

\fig{g06_p45_1}{Tiempos por vuelta (minutos) de los autos A, B y C: vuelta 1: 1,573 / 1,580 / 1,593; vuelta 2: 1,644 / 1,592 / 1,632; vuelta 3: 1,790 / 1,682 / 1,604.}

Los competidores ganan puntos de acuerdo con las siguientes reglas: Finalizada la segunda vuelta se dan 10 puntos de bonificación en la clasificación general a quien vaya de líder y 5 puntos a quien vaya de segundo. Al ganador de la competencia le dan 20 puntos, al segundo 10 puntos y al tercero 5 puntos.

a) Escribe cómo se distribuyen los puntos al finalizar la segunda vuelta y al finalizar la carrera. b) Determine cuál de los autos se acercó más a un minuto y sesenta décimas de minuto en la primera vuelta. c) Explica por qué considera que la práctica generalizada en este tipo de carreras el tiempo por vuelta se representa con tres dígitos después de la coma y no por dos. Calcula la diferencia de tiempo de los tres carros A, B y C en la primera vuelta con un carro D si se sabe que en esa vuelta invierte 1,4 minutos. ¿Cuáles serían estas diferencias si por un percance mecánico demora 1,09 minutos?

% PDF p. 45
\dba{2}{Utiliza las propiedades de los números enteros y racionales y las propiedades de sus operaciones para proponer estrategias y procedimientos de cálculo en la solución de problemas.}

\evidencias

\begin{itemize}
  \item Propone y utiliza diferentes procedimientos para realizar operaciones con números enteros y racionales.
  \item Argumenta de diversas maneras la necesidad de establecer relaciones y características en conjuntos de números (ser par, ser impar, ser primo, ser el doble de, el triple de, la mitad de, etc).
\end{itemize}

\ejemplo

Un obrero tiene que controlar la cantidad de vapor que hay en la caldera a través de un dispositivo conectado a ella, así como muestra la figura. Por este dispositivo entra vapor que se encarga de mover el corcho. El corcho sube cuando aumenta la cantidad de vapor en la caldera y baja cuando disminuye. Cuando el nivel del vapor en la caldera es el normal el corcho marca exactamente el punto cero.

\fig{g06_p46_1}{Caldera con un dispositivo medidor (corcho que sube y baja) y un obrero que la controla.}

\begin{itemize}
  \item Para indicar la ubicación del corcho se utilizan números acompañados del signo más (+) o del signo menos (-). Si el corcho está por encima del punto cero su ubicación se representa con un número acompañado del signo más. Si el corcho está por debajo de cero la ubicación del corcho se representa con un número acompañado del signo menos.
  \item Los desplazamientos del corcho hacia arriba se representan por números acompañados del signo “+” Los desplazamientos hacia abajo se representan por números acompañados del signo “-”. Calcula el desplazamiento del corcho si inicialmente está en la raya -23 y después aparece en la raya marcada con -45. Identifica cuál de las dos expresiones, (-80) + 8(-15) y (-80) - 8(-15), permite calcular la raya final a la que llega el corcho si inicialmente está en la raya marcada con -80 y a partir de este punto hace 8 desplazamientos uno tras otro hacia abajo. Cada desplazamiento tiene una magnitud de 15 rayas. Da razones de su elección 1 .
\end{itemize}
1 Adaptación de una tarea tomada de Pruebas Comprender de Matemática. Grado Noveno (2005) Secretaría de Educación de Bogotá D.C.

% PDF p. 46
\dba{3}{Reconoce y establece diferentes relaciones (orden y equivalencia) entre elementos de diversos dominios numéricos y los utiliza para argumentar procedimientos sencillos.}

\evidencias

\begin{itemize}
  \item Determina criterios de comparación para establecer relaciones de orden entre dos o más números.
  \item Representa en la recta numérica la posición de un número utilizando diferentes estrategias.
  \item Describe procedimientos para resolver ecuaciones lineales.
\end{itemize}

\ejemplo

En la figura se muestra una secuencia de imágenes que ilustran formas de encontrar el valor de x.

\fig{g06_p46_2}{Tres estados del dispositivo medidor con el corcho a distintas alturas.}

Describe diferentes procedimientos o acciones que le permitan conocer el valor de x y pone a prueba esos procedimientos.

% PDF p. 46
\dba{4}{Utiliza y explica diferentes estrategias (desarrollo de la forma o plantillas) e instrumentos (regla, compás o software) para la construcción de figuras planas y cuerpos.}

\evidencias

\begin{itemize}
  \item Construye plantillas para cuerpos geométricos dadas sus medidas.
  \item Selecciona las plantillas que genera cada cuerpo a partir del análisis de su forma, sus caras y sus vértices.
  \item Utiliza la regla no graduada y el compás para dibujar las plantillas de cuerpos geométricos cuando se tienen sus medidas.
\end{itemize}

\ejemplo

Se quieren forrar con papel de colores unos cuerpos geométricos como los que se muestran en la imagen.

\fig{g06_p47_1}{Dos estudiantes junto a cuerpos geométricos de colores (prismas, pirámides) y moldes planos, que se quieren forrar con papel.}

Selecciona de los moldes que se muestran en la figura los que se podrían utilizar para construir, con regla y compás, los forros respectivos y determina cómo calcular la cantidad de papel que se requiere para elaborar cada forro, si se conocen las medidas de las aristas de cada cuerpo.

Aunque estos tres cuerpos tuvieran igual medida en sus aristas, el volumen de los tres sería diferente. Encuentra el de mayor volumen y explica la respuesta.

% PDF p. 47
\dba{5}{Propone y desarrolla estrategias de estimación, medición y cálculo de diferentes cantidades (ángulos, longitudes, áreas, volúmenes, etc.) para resolver problemas.}

\evidencias

\begin{itemize}
  \item Decide acerca de las estrategias para determinar qué tan pertinente es la estimación y analiza las causas de error en procesos de medición y estimación.
  \item Estima el resultado de una medición sin realizarla, de acuerdo con un referente previo y aplica el proceso de estimación elegido y valora el resultado de acuerdo con los datos y contexto de un problema.
  \item Estima la medida de longitudes, áreas, volúmenes, masas, pesos y ángulos en presencia o no de los objetos y decide sobre la conveniencia de los instrumentos a utilizar, según las necesidades de la situación.
\end{itemize}

\ejemplo

Se presentan dos imágenes, una con un jugador frente al arco y la otra con 7 jugadores frente al arco para practicar sus lanzamientos. Estima la medida de los ángulos que forma cada jugador con respecto a los dos palos del arco y argumenta en qué posición existe mayor posibilidad de gol. Utiliza el transportador para medir los ángulos y compara esas medidas con las de la estimación, explica las estrategias utilizadas en ambos casos.

\fig{g06_p47_2}{Portería de fútbol vista desde arriba con los ángulos de tiro de dos jugadores.}

\fig{g06_p47_3}{Barrera de 7 jugadores numerados en fila frente a la portería.}

\fig{g06_p47_4}{Barrera de jugadores dispuesta en diagonal frente a la portería.}

% PDF p. 48
\dba{6}{Representa y construye formas bidimensionales y tridimensionales con el apoyo en instrumentos de medida apropiados.}

\evidencias

\begin{itemize}
  \item Diferencia las propiedades geométricas de las figuras y cuerpos geométricos.
  \item Identifica los elementos que componen las figuras y cuerpos geométricos.
  \item Describe las congruencias y semejanzas en figuras bidimensionales y tridimensionales.
  \item Estima áreas y volúmenes de figuras y cuerpos geométricos.
  \item Construye cuerpos geométricos con el apoyo de instrumentos de medida adecuados.
\end{itemize}

\ejemplo

Los productos de la industria son envasados en diferentes materiales: cartón, vidrio, plástico, metal y diferentes formas. A partir de las condiciones reales identifica las formas volumétricas que los constituyen, construye representaciones semejantes para configurar sus desarrollos geométricos y estima la cantidad de material necesario para su fabricación. Compara la información sobre volumen y peso que aparece en algunos empaques y establece relaciones entre ellas.

\fig{g06_p48_1}{Botellas y vasos con líquido de distintas capacidades.}

% PDF p. 48
\dba{7}{Reconoce el plano cartesiano como un sistema bidimensional que permite ubicar puntos como sistema de referencia gráfico o geográfico.}

\evidencias

\begin{itemize}
  \item Localiza, describe y representa la posición y la trayectoria de un objeto en un plano cartesiano.
  \item Identifica e interpreta la semejanza de dos figuras al realizar rotaciones, ampliaciones y reducciones de formas bidimensionales en el plano cartesiano.
\end{itemize}

\ejemplo

Elabora diseños de bisutería artesanal para crear diferentes pulseras con diversos materiales. Utiliza el plano cartesiano para identificar patrones y los expresa como parejas ordenadas y modifica estos patrones para producir nuevos modelos.

\fig{g06_p48_2}{Patrón de tejido en cuadrícula (mosaico de chaquiras) con simetrías.}

Tomado de: https://www.youtube.com/watch?v=IqVR8\_Tmjc.

% PDF p. 49
\dba{8}{Identifica y analiza propiedades de covariación directa e inversa entre variables, en contextos numéricos, geométricos y cotidianos y las representa mediante gráficas (cartesianas de puntos, continuas, formadas por segmentos, etc.).}

\evidencias

\begin{itemize}
  \item Propone patrones de comportamiento numéricos y expresa verbalmente o por escrito los procedimientos matemáticos.
  \item Realiza cálculos numéricos, organiza la información en tablas, elabora representaciones gráficas y las interpreta.
  \item Trabaja sobre números desconocidos y con esos números para dar respuestas a los problemas.
\end{itemize}

\ejemplo

Yadira se mudará a otro apartamento y pide cotizaciones a tres empresas de transporte. Cada empresa da sus condiciones:

Hasta 300 Hasta 200 Hasta 400 kilos para kilos para kilos para tarifa 300.000 tarifa de tarifa 400.000

\fig{g06_p49_2}{Tarifas de transporte: Empresa A hasta 300 kg \$300.000, más de 300 kg \$300.000 + \$2.000 por kilo extra; Empresa B hasta 200 kg \$250.000, más \$1.000 por kilo extra; Empresa C hasta 400 kg \$400.000, más \$500 por kilo extra.}

Más de 300 Más de 200 Más de 400 kilos: 300.000 kilos: 250.000 kilos: 600.000 pesos más pesos más pesos más 500 2.000 pesos 1.000 pesos pesos por cada por cada kilo por cada kilo kilo de más. de de

Representa en el plano cartesiano el costo total de transporte en pesos, en términos del peso transportado. Averigua en cuál de las empresas la razón entre peso transportado y costo es la mayor. Representa gráficamente las ofertas de las tres empresas en un diagrama cartesiano. Decide cuál es la empresa que le conviene contratar a Yadira, si tiene 400 kilogramos para transportar, o si el peso es de 700 kilogramos.

[1] Adaptado de matemáticas de la Vida Real (2011). G. Barozzi; M, Bergamini; D, Boni; R, Ceriani; L. Pagani. Editorial Octaedro.

% PDF p. 50
\dba{9}{Opera sobre números desconocidos y encuentra las operaciones apropiadas al contexto para resolver problemas.}

\evidencias

\begin{itemize}
  \item Utiliza las operaciones y sus inversas en problemas de cálculo numérico.
  \item Realiza cálculos numéricos, organiza la información en tablas, elabora representaciones gráficas y las interpreta.
  \item Realiza combinaciones de operaciones, encuentra propiedades y resuelve ecuaciones en donde están involucradas.
\end{itemize}

\ejemplo

Una compañía de pintura contrata empleados por días. La compañía determina que el monto que se paga por hora trabajada es de \$8.000, en jornada normal (8 horas diarias), pero si se hacen horas extras, se paga la hora a \$9.000 (máximo 4 diarias). Describe verbal, numérica, gráfica o simbólicamente, el monto que se ha de pagar diariamente y en varios días según la cantidad de horas extras. Calcula el monto de contribución a Seguridad Social en función del dinero ganado. La regla de la Oficina de Seguridad Social dice: Si se trabaja 8 horas diarias, se ha de pagar \$1.000 fijos, más un 5\% para salud, pensiones y cesantías. Pero si se trabajan más de 8 horas, se ha de pagar \$1.000, más un 4\% para salud, pensiones y cesantías.

Representa esos resultados en una tabla y una gráfica cartesiana, utiliza esta información para determinar el número de horas trabajadas por una persona que ganó en un día \$99.000. Si una persona debe pagar \$5.320 de aportes a la Seguridad Social, determina, cantidad de horas que trabajó y si trabajó horas extras.

% PDF p. 50
\dba{10}{Interpreta información estadística presentada en diversas fuentes de información, la analiza y la usa para plantear y resolver preguntas que sean de su interés.}

\evidencias

\begin{itemize}
  \item Lee y extrae la información estadística publicada en diversas fuentes.
  \item Plantea una pregunta que le facilite recolectar información que le permita contrastar la información estadística publicada.
  \item Organiza la información recolectada en tablas y la representa mediante gráficas adecuadas.
  \item Calcula las medidas requeridas de acuerdo a los datos recolectados y usa, cuando sea posible, calculadoras o software adecuado.
  \item Escribe un informe en el que analiza la información presentada en el medio de comunicación y la contrasta con la obtenida en su estudio.
\end{itemize}

\ejemplo

A finales del 2012 en su informe trimestral el Ministerio de Tecnologías de la Información y las comunicaciones, MINTIC, publicó la siguiente información:

“En cuanto a la telefonía móvil, sector que tuvo un crecimiento de 1,13\%, en Colombia hay más de 48 millones de abonados, lo que quiere decir que en el país hay 104,5 líneas por cada 100 habitantes. De estos 18,86\% se encuentra en la modalidad de postpago, y 81,14\% es prepago” (Tomado de http://colombiatic.mintic.gov.co/602/ articles-15179\_archivo\_pdf.pdf).

Responde preguntas como ¿el comportamiento del uso de la telefonía móvil en el salón es similar a lo afirmado en la noticia? Si se presentan diferencias con los datos de la noticia, ¿cuáles son las posibles razones para que esto suceda?

% PDF p. 51
\dba{11}{Compara características compartidas por dos o más poblaciones o características diferentes dentro de una misma población para lo cual seleccionan muestras, utiliza representaciones gráficas adecuadas y analiza los resultados obtenidos usando conjuntamente las medidas de tendencia central y el rango.}

\evidencias

\begin{itemize}
  \item Comprende la diferencia entre la muestra y la población.
  \item Selecciona y produce representaciones gráficas apropiadas al conjunto de datos, usando, cuando sea posible, calculadoras o software adecuado.
  \item Interpreta la información que se presenta en los gráficos usando las medidas de tendencia central y el rango.
  \item Compara las características de dos o más poblaciones o de dos o más grupos, haciendo uso conjunto de las respectivas medidas de tendencia central y el rango.
  \item Describe el comportamiento de las características de dos o más poblaciones o de dos o más grupos de una población, a partir de las respectivas medidas de tendencia central y el rango.
\end{itemize}

\ejemplo

En un estudio reciente sobre la velocidad de descarga de fotos y videos en celulares, se sometieron a prueba dos marcas, durante una hora se usaron 10 celulares diferentes de cada marca. Si se desea comprar un celular con una buena velocidad de descarga, ¿cuál de las dos marcas seleccionaría? Interpreta la información representada en los gráficos, y utiliza las medidas adecuadas para realizar la comparación que le permita tomar una buena decisión y así justificar las ventajas de una marca sobre la otra.

\fig{g06_p51_1}{Diagrama de barras horizontales: velocidad de descarga en Mbps (100–150) para 10 mediciones.}

\fig{g06_p52_1}{Segundo diagrama de barras horizontales de velocidad de descarga en Mbps para 10 mediciones.}

% PDF p. 52
\dba{12}{A partir de la información previamente obtenida en repeticiones de experimentos aleatorios sencillos, compara las frecuencias esperadas con las frecuencias observadas.}

\evidencias

\begin{itemize}
  \item Enumera los posibles resultados de un experimento aleatorio sencillo.
  \item Realiza repeticiones del experimento aleatorio sencillo y registra los resultados en tablas y gráficos de frecuencia.
  \item Interpreta y asigna la probabilidad de ocurrencia de un evento dado, teniendo en cuenta el número de veces que ocurre el evento en relación con el número total de veces que realiza el experimento.
  \item Compara los resultados obtenidos experimentalmente con las predicciones anticipadas.
\end{itemize}

\ejemplo

En un el juego con dados, uno en forma de cubo (6 caras) y otro en forma de tetraedro (4 caras) participan dos jugadores. Cada uno selecciona un dado, lo lanza al aire y gana quien obtenga más veces el número 1, después de hacer 100 lanzamientos. Leidy dice que ella juega si selecciona el dado en forma de tetraedro. Justifica la selección de uno de los dados, anticipa la posibilidad de ocurrencia del evento que salga el 1, realiza el experimento, registra los resultados y compara y razona sobre las diferencias entre lo esperado y lo observado. Verifica la validez de la afirmación de Leidy.

